%\usepackage{pslatex} % to see correctly a .pdf file on the computer screen
\usepackage{pgf}
\usepackage{color}
\usepackage{graphicx}
\usepackage{amssymb} %symbole de maths
\usepackage{amsmath} %idem
\usepackage{upgreek} %idem
\usepackage[utf8]{inputenc}
\usepackage{listings} % code display
\usepackage{marvosym} %\MVRightarrow
%\usepackage{fancyvrb} %give size to verbatim
%\usepackage{hyperref}
\usepackage[english,francais]{babel}
\definecolor{vertmoyen}{RGB}{51,110,23} % vert moyen
\usecolortheme[named=vertmoyen]{structure}
\usepackage{tabularx} % varier la largeur du tableau
\usepackage{layout}
\usepackage{longtable}
\setlength{\LTleft}{-5cm plus 1 fill}
\setlength{\LTright}{-5cm plus 1 fill}
\usepackage{booktabs}
\usepackage{arydshln} %% dashlines for tabular
\newcommand{\logit}{\text{logit}}
\newcommand{\bs}[1]{\boldsymbol{#1}}
\newcommand{\R}{\textnormal{\sffamily\bfseries R}}
\newcommand{\pkg}[1]{{\fontseries{b}\selectfont #1}}
\newcolumntype{C}[1]{>{\centering\arraybackslash}m{#1}}
\newcommand{\red}[1]{\textcolor{red}{#1}}
\newcommand{\blue}[1]{\textcolor{blue}{#1}}
%% Natbib is a popular style for formatting references.
%\usepackage{natbib} %doesn't work with beamer

\title[jSDM R package]{\textbf{jSDM R package for Joint Species Distribution Models}}
%\subtitle{}
\date{}

% Theme
% \usetheme{AnnArbor}
% \usetheme{Dresden}
% \usetheme{Copenhagen}
% \usetheme{Frankfurt}
% \usetheme{Berlin}
% \usetheme{Madrid}
% \usetheme{Montpellier}
% \usetheme{Singapore}
% \usetheme{Antibes}
\useinnertheme{rounded} %% bullets
\useoutertheme[subsection=false]{miniframes}
%\useoutertheme[subsection=false,footline=authorinstitutetitle]{miniframes}
\setbeamertemplate{footline}[frame number]

% % Set beamer colors
% \definecolor{Vert}{RGB}{51,110,23} % UBC Blue (primary)
% \definecolor{Grey}{rgb}{0.3686, 0.5255, 0.6235} % UBC Grey (secondary)
% \setbeamercolor{palette primary}{bg=Vert,fg=white}
% \setbeamercolor{palette secondary}{bg=Vert,fg=white}
% \setbeamercolor{palette tertiary}{bg=Vert,fg=white}
% \setbeamercolor{palette quaternary}{bg=Vert,fg=white}
% \setbeamercolor{structure}{fg=Vert} % itemize, enumerate, etc
% \setbeamercolor{section in toc}{fg=Vert} % TOC sections
% 
% % Override palette coloring with secondary
% \setbeamercolor{subsection in head/foot}{bg=Grey,fg=white}

% Ignore ignorenonframetext class option in default template
\makeatletter
\beamer@ignorenonframefalse
\makeatother

% Logo
\newif\ifplacelogo % create a new conditional
\logo{
  \ifplacelogo
  \includegraphics[height=0.7cm]{figs/logo_Cirad.png}\hspace{0.25cm}
  \includegraphics[height=1cm]{figs/logo_AMAP.png}
  \fi
} % replace with your own command

%Call table of contents at the beginning of each section
\AtBeginSection[]{
\placelogotrue
  \begin{frame}
    \frametitle{Plan}
    \begin{columns}[c]
      \begin{column}{0.5\textwidth}
        \tableofcontents[sections=1, currentsection]
        \vspace{0.5cm}
        \tableofcontents[sections=2, currentsection]
      \end{column}
      \begin{column}{0.5\textwidth}
        \tableofcontents[sections=3, currentsection]
        \vspace{0.5cm}
        \tableofcontents[sections=4, currentsection]
      \end{column}
    \end{columns}
  \end{frame}
\placelogofalse
}

\AtBeginSubsection[]{}

% \AtBeginSubsection[]{
% \placelogotrue
%   \begin{frame}
%     \frametitle{Plan}
%     \begin{columns}[c]
%       \begin{column}{0.5\textwidth}
%         \tableofcontents[sections={1-2},currentsection,currentsubsection]
%       \end{column}
%       \begin{column}{0.5\textwidth}
%         \tableofcontents[sections={3-4},currentsection,currentsubsection]
%       \end{column}
%     \end{columns}
%   \end{frame}
% \placelogofalse
% }

% Two-columns slides
\def\bcols{\begin{columns}}
\def\bcol{\begin{column}}
\def\ecol{\end{column}}
\def\ecols{\end{columns}}

% Change fontsize of R code
\let\oldShaded\Shaded
\let\endoldShaded\endShaded
\renewenvironment{Shaded}{\footnotesize\oldShaded}{\endoldShaded}

% Change fontsize of output
\let\oldverbatim\verbatim
\let\endoldverbatim\endverbatim
\renewenvironment{verbatim}{\footnotesize\oldverbatim}{\endoldverbatim}
